\section{引言}

ProteinMPNN 以给定蛋白骨架为条件生成氨基酸序列，原始工作表明其能够处理单链与多链设计问题\cite{dauparas2022proteinmpnn}。本项目在此基础上增加了 N-甲基化表示：先预测天然氨基酸，再由与天然残基类别相对应的甲基化专家头给出位点置信度。因此，模型实际上同时面对两个相互耦合但不等价的问题：

\begin{enumerate}
  \item \textbf{基础序列问题：}给定结构，天然化后的氨基酸是否恢复正确；
  \item \textbf{甲基化问题：}在已知或模型预测的天然氨基酸条件下，该位点是否应标记为 N-甲基化；
  \item \textbf{结构问题：}生成序列形成的结构是否保持合理的内部构象，并处于正确的受体结合姿态。
\end{enumerate}

这三类问题不能用一个数字代替。较高的甲基分类 AUC 不等于较高的天然序列恢复率；较低的肽自对齐 RMSD 也不等于正确结合；结构预测文件中的 B-factor/pLDDT 代理值更不是物理能量。为了避免这些概念在先前零散汇报中互相混淆，本研究将所有可追溯证据按数据队列和评价定义重新组织。

本稿只完成 \taskone，研究对象严格冻结为原始 \path{frankenstein_v28.pt}。后续 V7、V8、V9 的训练、拼接或稳定性修复不属于本文结果。核心问题包括：

\begin{itemize}
  \item 原始模型在单体留出集上的 RAA、AUC、Accuracy、Precision、Recall 与 F1 到底是多少；
  \item 为什么先前截图出现 16.59\%、38.11\% 与 85.67\% 等值，以及它们能否用于论文；
  \item 在温度 0.5、甲基化置信度严格 $>0.6$ 的条件下，六个非 3AV 复合物从原始序列到联合 RMSD $<3$~\AA 或 $<5$~\AA 的真实产率是多少；
  \item 17 靶点的 best85 结构、置信度、多样性、通透性与 fixed-pose 能量能说明什么，又不能说明什么；
  \item 尚哥提出的全部指标中，哪些已经有合格证据，哪些仍未计算或在数学上不可定义。
\end{itemize}

论文的主张边界是保守且明确的：我们报告已经观测到的模型输出和结构计算结果，但不从预测分数外推实验透膜性、结合自由能或生物活性；不把缺失值填成零；不以选择后的代表样本替代完整生成分母。
